% problem-000091 7 含磷小分子反应
\section*{7. 含磷小分子反应}
作为⼈⼯合成的⼤环配体之⼀，咔咯的合成及其与⾦属离⼦的配位化学备受关注。利⽤咔咯作为反应配体，可实现过渡⾦属离⼦的⾮常规氧化态，制备功能材料，探索新型催化剂，等等。下图(a)给出三（五氟苯基）咔咯分⼦(A)的⽰意图，简写为 $\mathrm { H } _ { 3 } ( \mathrm { t p f c } )$ 。它与⾦属离⼦结合时四个氮原⼦均参与配位。室温下，空⽓中，$\mathrm { C r } ( \mathrm { C O } ) _ { \mathfrak { c } }$ 和A在甲苯中回流得到深红⾊晶体B（反应1），B显顺磁性，有效磁矩为 $1 . 7 2 \ \mu _ { \mathrm { B } } ,$ ，其中⾦属离⼦的配位⼏何为四⽅锥；在惰性⽓氛保护下，B与三苯基膦 $\left( \mathrm { P P h } _ { 3 } \right)$ 和三苯基氧膦 $\mathrm { ( O P P h _ { 3 } ) }$ 按1:1:1在甲苯中反应得到绿⾊晶体C（反应2）；在氩⽓保护下， $\mathrm { C r C l _ { 2 } }$ 和A在吡啶（简写为 $\mathrm { p y } )$ ）中反应，得到深绿⾊晶体D（反应3），D中⾦属离⼦为⼋⾯体配位，配位原⼦均为氮原⼦。在⼀定的条件下，B、C和D之间可以发⽣转化（下图b），这⼀过程被认为有可能⽤于 $\mathrm { O _ { 2 } }$ 的活化或消除。

（图：）
(a)

（图：）
(b)

\subsection*{7-1}
A中的咔咯环是否有芳⾹性？与之对应的π电⼦数是多少？

\subsection*{7-2}
写出B、C、D的化学式（要求：配体均⽤简写符号表⽰）。

\subsection*{7-3}
写出B、C、D中⾦属离⼦的价电⼦组态（均⽤ $d ^ { n }$ ⽅式表⽰）。

\subsection*{7-4}
写出制备B、C、D的反应⽅程式（要求：配体均⽤简写符号，系数为最简整数⽐）。

\subsection*{7-5}
利⽤电化学处理，B可以得电⼦转化为B−，也可以失去电⼦转化 $\mathbf { B } _ { \mathrm { ~ c ~ } } ^ { + }$ 。与预期的磁性相反，B+依然显⽰顺磁性。进⼀步光谱分析发现，与咔咯环配体相关的吸收峰位置发⽣显著了变化。推测⾦属离⼦的价电⼦组态（⽤ $\lceil d ^ { n }$ ⽅式表⽰），指出磁性与光谱变化的原因。
